% Generated draft: requires mathematical and source review.
\documentclass[11pt]{article}
\usepackage[margin=25mm]{geometry}
\usepackage{fontspec}
\setmainfont{Latin Modern Roman}
\usepackage{amsmath,amssymb,mathtools,unicode-math}
\setmathfont{Latin Modern Math}
\usepackage{enumitem,needspace,xurl,hyperref}
\setlength{\parindent}{0pt}
\setlength{\parskip}{5pt}
\setlength{\emergencystretch}{3em}
\newcommand{\N}{\mathbb{N}}
\newcommand{\Z}{\mathbb{Z}}
\newcommand{\Q}{\mathbb{Q}}
\newcommand{\R}{\mathbb{R}}
\newcommand{\C}{\mathbb{C}}
\newcommand{\F}{\mathbb{F}}
\newcommand{\A}{\mathbb{A}}
\newcommand{\PP}{\mathbb{P}}
\newcommand{\E}{\mathbb{E}}
\newcommand{\eps}{\varepsilon}
\newcommand{\dd}{\,\mathrm{d}}
\begin{document}
\section*{UnsolvedMath ID: 800405}
\textbf{Source problem code:} catalogue entry
\subsection{Definitions and mathematical statement}
Let \(X\) be a smooth projective variety of dimension \(N\), and let \(f\colon X\dashrightarrow X\) be a dominant rational map. The supplied statement does not specify the ground field; its intended field hypotheses remain unresolved.

For the usual intersection-theoretic meaning of the higher dynamical degrees, choose an ample Cartier divisor \(H\) on \(X\). For each integer \(n\geq 1\), let \(\Gamma_n\subseteq X\times X\) be the closure of the graph of the iterate \(f^n\), and let \(p_{1,n},p_{2,n}\colon\Gamma_n\to X\) be the two projections. For \(1\leq i\leq N\), put
\[
 d_i(f^n;H)=\deg\!\left((p_{1,n}^{*}H)^{N-i}(p_{2,n}^{*}H)^i\cap[\Gamma_n]\right).
\]
The standard dynamical-degree convention is
\[
 \delta_f^{(i)}=\lim_{n\to\infty}d_i(f^n;H)^{1/n},
\]
in a setting where the standard existence of this limit and its independence of \(H\) apply. This definition uses the graph of each iterate separately. The supplied record does not itself define the dynamical degrees or specify the field hypotheses for this convention.

A number \(\alpha\) is an algebraic integer if there exist an integer \(m\geq 1\) and integers \(a_0,\ldots,a_{m-1}\) such that
\[
 \alpha^m+a_{m-1}\alpha^{m-1}+\cdots+a_1\alpha+a_0=0.
\]
\textbf{Question 28.} For every \(X\) and \(f\) as above, is \(\delta_f^{(i)}\) an algebraic integer for every integer \(i\) with \(1\leq i\leq N\)? The intended scope over ground fields must be clarified before this question has a fully specified domain.
\subsection{Short English statement}
Must every higher dynamical degree of a dominant rational self-map of a smooth projective variety be an algebraic integer? The supplied record labels this question open, but that label is not independently verified here. The intended ground field is unspecified.
\subsection{Sources}
\begin{itemize}
\item Local UnsolvedMath registry, permanent ID \(800405\), source problem code \texttt{catalogue entry}, titled \emph{Open Problems in Arithmetic Dynamics --- Question 28}. The supplied record attributes the question to Benedetto, DeMarco, Ingram, Jones et al., with proposed year 2018, and labels it open.
\item Benedetto, DeMarco, Ingram, Jones et al., \emph{Current Trends and Open Problems in Arithmetic Dynamics} (2018), Question 28. Supplied source URL: \url{https://arxiv.org/abs/1806.04980}. The source record reports an access date of 2026-10-16 and extraction method \texttt{arxiv\_latex}. These are supplied provenance details, not a claim of independent source access or verification.
\end{itemize}
\end{document}
